\subsection{结构群的扩张与限制}

6.6.1. 设 \(\lambda = (\mathbf{P}, \mathbf{G}, \mathbf{B}, \pi)\) 是主纤维化，且 \(\varphi\) 是从 \(\mathbf{G}\) 到群流形 \(\mathbf{H}\) 的同态。令 \(\mathbf{G}\) 按 \(\mathbf{H}\) 在 \(g \cdot h = \varphi(g) \cdot h\) 上左作用，并令 \(\mathbf{P} \times^{\mathbf{G}} \mathbf{H}\) 为与 \(\lambda\) 相关且典型纤维为 \(\mathbf{H}\) 的纤维丛。由于 \(\mathbf{H}\) 中的右平移与 \(\mathbf{G}\) 的运算相容，群 \(\mathbf{H}\) 在 \(\mathbf{P} \times^{\mathbf{G}} \mathbf{H}\) 上右作用（参见第 6.5.5 节）；若 \(\pi_{\mathbf{H}}\) 表示从 \(\mathbf{P} \times^{\mathbf{G}} \mathbf{H}\) 到 \(\mathbf{B}\) 的投影，则四元组（\(\mathbf{P} \times^{\mathbf{G}} \mathbf{H}\)，\(\mathbf{H}\)，\(\mathbf{B}\)，\(\pi_{\mathbf{H}}\)）是一个主纤维化，记为 \(\varphi(\lambda)\)；称其为通过同态 \(\lambda\) 由 \(\varphi\) 得到的。

把  对应到  的等价类的映射 \(f\) 是从 \(\mathbf{P}\) 到 \(\mathbf{P} \times^{\mathbf{G}}\mathbf{H}\) 的与 \(x \in \mathbf{P}\)\((x, e)\) 相容的 \(\mathbf{B}\)-态射（参见第 6.3.1 节）。此外，若  是从 \(\mathbf{P}\) 到结构群为 \(\mathbf{P} \times^{\mathbf{G}}\mathbf{H}\)\(\varphi\)\(f'\) 的主纤维丛  的 \(\mathbf{B}\)-态射，且 \(\mathbf{P}\)\(\mathbf{P}'\)\(\mathbf{H}\)\(f'\) 与 \(\varphi\) 相容，则存在唯一的从  到  的 \(\mathbf{H}\)-\(\mathbf{B}\)-同构 \(\theta\)，使得 \(\mathbf{P} \times^{\mathbf{G}}\mathbf{H}\)\(\mathbf{P}'\)\(f' = \theta \circ f\)。

6.6.2. 假设 \(\lambda\) 由 B 的开覆盖 \(\mathcal{V} = (\mathrm{U}_i)\) 和上循环 \((g_{ij})\)（6.4.2）定义。则 \(\varphi(\lambda)\) 可用同一覆盖和上循环 \((h_{ij})\) 定义，其中 \(h_{ij} = \varphi \circ g_{ij}\)。

6.6.3. 设 F 是一个流形，群 H 在其上左作用；用 \((h, y) \mapsto h \cdot y\) 表示 H 在 F 上的作用律。群 G 通过 \((g, y) \mapsto \varphi(g) \cdot y\) 在 F 上作用。设 E 是与 \(\varphi(\lambda)\) 相关且典型纤维为 F 的纤维丛。映射 \((x, y) \mapsto f(x) \cdot y\) 从 P × F 到 E（映射 f 即第 6.6.1 节中定义的映射）赋予 E 一个与 \(\lambda\) 相关且纤维类型为 F 的纤维丛结构。特别地，\((\mathbf{P} \times^{\mathbf{G}}\mathbf{H}) \times^{\mathbf{H}}\mathbf{F}\) 可与 P × \(^{\mathbf{G}}\) F 识别。

6.6.4. 假设 G 是 H 的李子群，且 \(\varphi: \mathbf{G} \to \mathbf{H}\) 是 G 到 H 的典范嵌入。此时称 \(\varphi(\lambda)\) 是通过将结构群扩张为 H 由 \(\lambda\) 得到的。第 6.6.1 节的态射 \(f: \mathbf{P} \to \mathbf{P} \times^{\mathbf{G}} \mathbf{H}\) 是从 P 到 \(\mathbf{P} \times^{\mathbf{G}} \mathbf{H}\) 的一个同构，其像为后者的闭子流形（当 \(\mathbf{P} \times^{\mathbf{G}} \mathbf{H}\) 时等于 \(\mathbf{G} = \mathbf{H}\)）；这个同构与 G 的运算相容（注意 G 作为 H 的子群在 \(\mathbf{P} \times^{\mathbf{G}} \mathbf{H}\) 上作用）。

6.6.5. 进一步假设 G 是 H 的李子群，并设 \(\mu = (Q, H, B, \pi)\) 是结构群为 H、基空间为 B 的主纤维化，E 是与 \(\mu\) 相关且纤维类型为 H/G 的纤维丛。若 \(\gamma\) 表示从 H 到 H/G 的典范投影，则对  置 \(\delta(y) = y \cdot \gamma(e)\)（其中 e 表示 H 的单位元）；于是得到一个态射 \(y \in Q\)\(e\)\(\delta: Q \to E\)，且四元组 \((Q, G, E, \delta)\) 是一个主纤维化。特别地，Q/G 可典范地识别为 E，而 E 又同构于 \(Q \times^{H} H/G^{1}\) 。

现在设 \(s: \mathbf{B} \to \mathbf{E}\) 是 \(\mathbf{E}\) 的一个截面，并设 \(\lambda_s\) 是我们刚刚定义的纤维化 \(s\) 在 \((\mathbf{Q}, \mathbf{G}, \mathbf{E}, \delta)\) 下的逆像。它是一个以 \(\mathbf{B}\) 为基空间、结构群为 \(\mathbf{G}\) 的主纤维化，并且其向 \(\mathbf{H}\) 的扩张同构于 \(\mu\)；具有这些性质的任意主纤维化都可以通过这种方式在同构意义下得到；\(s_1\) 的两个截面 \(s_2\) 和 \(\mathbf{E}\) 定义同构的纤维化，当且仅当它们可由 \(\mu\) 的一个 H-B-自同构相互变换。

